% =====================================================================
%  CH. 7 LES CÉRÉMONIES · CH. 8 LE RUNBOOK DU MAÎTRE
% =====================================================================
\section{Les cérémonies — le maître fait, l'outil relaie}\label{sec:ceremonies}

\subsection{Le principe : personne ne touche vos secrets à votre place}

Toutes les cérémonies du coffre obéissent à la même discipline (D105) : le bac à
sable — l'agent, l'assistant, whatever — ne touchera \emph{jamais} le conteneur ni
les phrases réels. Les cérémonies sont des actes du maître, guidés par un outillage
qui \emph{relaie} sans réimplémenter : le GUI lance les vérités vivantes
(\code{coffre\_cf2.py}, \code{coffre\_cf3.py}, \code{sceau\_I.py}) et leur sortie
coule à l'écran. La phrase passe par STDIN — jamais en ligne de commande (invisible
des journaux de processus), jamais affichée, jamais dans un rapport. Le verdict est
\textbf{croisé} : code de sortie \emph{et} marque attendue dans le flux — un enfant
menteux qui affiche VERT mais sort 1 est ROUGE ; un enfant muet est SANS VERDICT.

Chaque module à tester porte le \textbf{trio D122} : une GUI, une boîte noire locale
(journal horodaté, gitignoré, muet si le disque refuse — témoin, pas dépendance), un
bouton « Envoyer le rapport ». Le rapport emporte l'écran et la boîte noire, aucune
phrase, aucun chemin de coffre réel.

\subsection{La cérémonie du sceau de l'Intention (quatre gestes)}\label{sec:cersceau}

Le sceau signe la spécification elle-même — l'arbre \code{v06/I/} et les treize
arbres de lignée : le garde est à QUATORZE arbres signés (D177 — l'arbre
\code{loupe} entré à la table canonique par la cérémonie d'ère du TGN).
Quatre gestes, depuis le clone :

\begin{enumerate}
\item \textbf{Préparer} (le dépôt) : \code{sceau\_I.py preparer} reconstruit le
manifeste de chaque arbre, compare au manifeste signé, écrit \code{a-signer.txt} —
une ligne par arbre : arbre, balise, empreinte.
\item \textbf{Signer} (hors ligne, votre clé) : la cérémonie autoportée
\code{ceremonie\_hors\_ligne.py signer} lit \code{a-signer.txt}, signe chaque
empreinte avec la clé privée physique, \emph{vérifie chaque signature avant
d'écrire} \code{signatures-I.txt}.
\item \textbf{Déposer} (au dépôt) : \code{sceau\_I.py deposer} installe la remise —
chaque signature est vérifiée contre l'arbre vivant et la clé publique AVANT toute
écriture ; une remise invérifiable n'écrit RIEN.
\item \textbf{Vérifier} : \code{sceau\_I.py verifier} — dix arbres VERTS, code 0.
Puis commit + push (la preuve est publique).
\end{enumerate}

Le dépôt ne voit \emph{jamais} la graine : le flux est air-gap, la charge utile est
imprimable (32 octets hex par arbre), la clé ne voyage pas. L'horodatage de première
intention : la date du commit signé.

\subsection{La cérémonie CF1 $\to$ CF2}\label{sec:cercf2}

Depuis le GUI de cérémonie (\code{lancer\_ceremonie}), onglet « Coffre CF2 » — le
coffre réel CF1 étant posé au bon endroit (le \code{.keys\textbackslash} du clone, ou
le dossier de cérémonie) :

\begin{enumerate}
\item \textbf{La phrase d'abord.} Le fichier-remise D81 est à jour (si vous venez de
changer la phrase : le fichier AVANT tout — un premier \code{migrer} refusé « phrase
de contrôle divergente » est le fail-closed qui travaille, pas une panne).
\item \textbf{Migrer.} Le CF1 est déverrouillé une dernière fois ; la base est
re-scellée sous M + Shamir. Naissent \code{coffre\_cf2.bin} et les cinq gardes
\code{g1} à \code{g5} (sous \code{gardes-cf2/}), tout dans \code{.keys}. Le CF1
reste \textbf{intact}.
\item \textbf{Trois copies hors-ligne.} Export de la feuille (fragments g3--g5, tags
vérifiés avant impression) : papier au sachet, et/ou fichiers \code{.cf2f} copiés sur
une USB hors-ligne. \emph{Trois}, pas deux : la relève exige le seuil.
\item \textbf{Contrôle de relecture.} Relisez les copies (le GUI guide) — une copie
illisible aujourd'hui est une résurrection impossible demain.
\item \textbf{Liste blanche} (sentinelle CF2) puis \textbf{CF2 — verifier} : VERT.
Sur Windows, la sentinelle refuse honnêtement de démarrer (moteur SACL/ETW restant) :
la porte D120 avertit « SANS sentinelle » à chaque déverrouillage — borne dite,
attente assumée.
\end{enumerate}

\begin{terrain}[— un refus nommé, ça ressemble à quoi]
Le premier \code{migrer} CF2 du coffre réel (31/10) : « ÉCHEC — phrase de contrôle
divergente ». La remise D81 n'était pas encore remplacée. Le refus est NOMMÉ, rien
n'est écrit, le coffre est intact — la cérémonie se rejoue après le geste de
discipline. C'est le fail-closed : \emph{mieux vaut un refus qu'un scellé sous une
phrase qui n'est pas celle que vous pensiez taper}.
\end{terrain}

\subsection{La cérémonie CF2 $\to$ CF3 — MIGRER, VÉRIFIER, DÉTRUIRE}\label{sec:cercf3}

Trois onglets, dans cet ordre, depuis \code{lancer\_migration\_cf3} (à la racine du
clone ou du kit) :

\begin{enumerate}
\item \textbf{MIGRER (actes 1--3).} Champs : le conteneur CF2 (pré-rempli
\code{.keys\textbackslash{} \allowbreak coffre\_cf2.bin}), la racine des gardes (vide = la
config), la config explicite (vide = \code{.keys\textbackslash config\_cf2.json} si
elle vit), le \textbf{média du piège — REQUIS} (branchez la clef USB du sachet et
désignez-la : la machine ne le crée jamais, spec § 4), le conteneur CF3 \emph{qui
naît} (dialogue d'enregistrement — c'est une sortie, pas une entrée), et la phrase
\textbf{deux fois} (un scellé sous une phrase tapée ne s'ouvrirait plus sous la
phrase pensée). La base re-scellée est de schéma identique — rien à traduire ; le
piège est posé sur le média ; la re-lecture à froide du staging doit être VERTE ; le
CF2 reste \textbf{intact}.
\item \textbf{VÉRIFIER (verdict à froid).} Le conteneur CF3, le média, la phrase —
une fois. Le verdict affiche époque, checksum, projet(s) : \textbf{notez le checksum
sur la carte du sachet}. Recommandé (D127) : re-exécutez le lendemain, machine
froide — un état qui ne se vérifie qu'à chaud n'est pas un état reposé.
\item \textbf{DÉTRUIRE (acte 4, explicite).} Uniquement après recoupement : le
conteneur CF2 à détruire, le conteneur CF3 \emph{survivant} (garde fail-closed : refus
si le remplaçant est absent), la case « Je comprends ». Les fragments CF2 et leurs
copies hors-ligne ne sont \emph{jamais} touchés — ils restent la relève historique.
\end{enumerate}

\begin{terrain}[— l'habit apprend du terrain : la grille D132]
Le GUI de migration v1.0 packait les champs pleine largeur puis les boutons
« Parcourir\ldots{} » à droite — la cavité les empilait en colonne, aucun sur la
ligne de son champ. Le retour du terrain (31/10) a enfanté la v1.1 : chaque ligne
de saisie en grille, le champ et son bouton sur la MÊME ligne, tout champ de chemin
porte son bouton, la confirmation de phrase est étiquetée, les boutons d'acte sont
pleine largeur. Vérifiable à l'ouverture : la fenêtre s'annonce
« migration-CF3-GUI v1.1 (D132) » — barre de titre, bandeau, et première ligne de la
boîte noire.
\end{terrain}

\subsection{La clôture, cérémonie du quotidien}\label{sec:cloture}

L'usage courant d'un coffre CF3 : \code{ouvrir} (média branché, phrase), travailler,
\code{fermer} (média rebranché, phrase retapée). Chaque clôture : une ligne sur la
carte (checksum), un piège de plus au sachet, une époque de plus au compteur. Entre
les deux, le brouillon vit sous K\_session — une panne le rend déchet, le canonique
reste scellé. La règle d'hygiène : sessions courtes, re-verrouillage rapide, média
débranché pendant le travail.

\section{Le runbook du maître — Windows, pas à pas}\label{sec:runbook}

\subsection{Trouver la maison}\label{sec:maison}

Toute cérémonie vit dans une \textbf{maison} : le clone git du dépôt (celui qui
contient \code{v06\textbackslash}) ou le kit de cérémonie du Drive. Le signe de la
maison : les lanceurs \code{.bat} à la racine, et \code{v06\textbackslash} plein
(42 fichiers, \code{I\textbackslash} et \code{sceau-I\textbackslash} compris). Le
lanceur cherche \code{v06\textbackslash migration\_cf3\_gui.py} \emph{d'abord}, puis
une copie plate à côté de lui — lancé depuis une maison, il prend la version vivante ;
lancé depuis un dossier où le \code{.py} a été posé seul, il vivrait en repli.

\begin{terrain}[— le GUI seul dans un dossier n'est pas une maison]
Le rapport du 30/10 20:03 : « modules REQUIS absents : coffre\_cf3.py,
coffre\_cf2.py », dossier v06 = \code{\ldots{}\textbackslash GitHub\textbackslash
Cérémonie}. Le GUI v1.1 avait été copié seul dans un dossier : le relais ne trouvait
pas ses vérités vivantes — refus nommé à chaque acte, par construction. Le remède :
lancer depuis la maison (le clone \code{triade}, ou le kit), où
\code{v06\textbackslash} est entier. Et l'autre doute du même soir — « le .bat
n'appelle pas le .py mis à jour » — se lît dans le lanceur : il appelle
\code{v06\textbackslash migration\_cf3\_gui.py} en premier ; le correctif v1.1 vivait
dans ce fichier même. La preuve à l'exécution : la fenêtre s'annonce v1.1 (D132),
trois fois.
\end{terrain}

\begin{table}[htbp]
\centering\small
\caption{Deux maisons, deux postures — le clone git et le kit Drive.}
\label{tab:maisons}
\begin{tabularx}{\textwidth}{@{}l Y Y@{}}
\toprule
 & \textbf{Le clone git} & \textbf{Le kit Drive} \\
\midrule
Qu'est-ce & le dépôt complet (\code{git clone}), \code{v06\textbackslash} entier à
jour par \code{git pull} & le dossier \code{/Cérémonie} du Drive : lanceurs,
LISEZMOI, \code{v06\textbackslash} identique au miroir PAR CONSTRUCTION (D130) \\
Se met à jour par & \code{git pull} (réflexe avant toute cérémonie) & le hook
post-commit du dépôt (miroir PUIS kit, même geste) — re-télécharger en cas de doute \\
Sert à & les cérémonies ET les remises du sceau (commit, push) & les cérémonies
seules — aucune écriture au dépôt \\
Comment savoir qu'il est frais & \code{git log} ; la version que le GUI s'annonce &
\code{MANIFESTE.txt} : commit de référence + empreintes — un kit périmé se LIT, ne
se devine pas \\
\bottomrule
\end{tabularx}
\end{table}

\subsection{L'ordre des gestes, du début à la fin}\label{sec:ordre}

Pour un coffre qui démarre de zéro (la chaîne complète, telle que le terrain l'a
parcourue) :

\begin{enumerate}
\item \textbf{Préparer la maison} : \code{git pull} dans le clone (ou re-télécharger
le kit) ; vérifier le \code{MANIFESTE.txt} du kit (commit de référence, empreintes)
si vous partez du Drive.
\item \textbf{Poser le CF1} : \code{coffre\_fort.bin} (le conteneur réel, livré à la
racine du coffre Drive) dans le \code{.keys\textbackslash} du clone. CF1 — statut
VERT.
\item \textbf{La phrase à vous} (optionnel, D110) : changer la phrase de contrôle
pour la vôtre — alphanumérique, sans plafond. Le fichier-remise D81 suit
\emph{immédiatement} : une ligne dans \code{PHRASE-CONTROLE-\allowbreak COFFRE.txt}, commit, push.
\item \textbf{CF1 $\to$ CF2} : cérémonie § \ref{sec:cercf2} — migrer, trois copies
hors-ligne, relecture, liste blanche, verifier VERT.
\item \textbf{CF2 $\to$ CF3} : cérémonie § \ref{sec:cercf3} — MIGRER (média du piège
requis, phrase deux fois), VÉRIFIER (checksum sur la carte), et le lendemain machine
froide : re-VÉRIFIER, recouper œil contre carte.
\item \textbf{DÉTRUIRE CF2} : après recoupement — la case explicite. Le rapport
part par « Envoyer le rapport ».
\end{enumerate}

Pour un \textbf{coffre neuf} : le niveau est \textbf{votre arbitrage à la création}
(D137, spécification § 12.5) — CF1 (\code{coffre\_fort.py init}, la relève simple),
CF2 (\code{init} CF1 puis la cérémonie \code{coffre\_cf2.py migrer}, la garde
répartie), CF3 (\code{coffre\_cf3.py init} : époque 1, média du piège, première
ligne de la carte). La chaîne ci-dessus est l'histoire du coffre réel, pas un rite
obligé — et elle reste disponible si les besoins grandissent.

\subsection{La check-list pré-vol — avant toute cérémonie}\label{sec:preflight}

Cinq vérifications, dans l'ordre, cinq minutes :

\begin{enumerate}
\item \textbf{La maison est à jour} : \code{git pull} rend « Already up to date »
(le GUI s'annonce avec sa version — v2.2 pour la cérémonie, v1.1 pour la migration) ;
le kit du Drive porte son \code{MANIFESTE.txt} au commit de référence.
\item \textbf{Le fichier-remise D81 dit la vérité} : la ligne de
\code{PHRASE-CONTROLE-COFFRE.txt} est la phrase que vous allez taper — sinon, posez-la
d'abord, commit, push.
\item \textbf{Le sachet est à portée de main} : média du piège à brancher au bon
moment (et à débrancher après), carte pour noter le checksum, fragments hors-ligne si
la cérémonie les touche.
\item \textbf{La pile est verte} : le rapport d'environnement du GUI dit
« pile COMPLÈTE » (argon2-cffi, PyNaCl, cryptography) — un module manquant est un
refus nommé à l'acte, autant le savoir avant.
\item \textbf{Vous avez le temps} : les cérémonies ne se pressent pas ; la phrase se
tape deux fois précisément PARCE QUE le scellé sous une phrase mal tapée ne
s'ouvrirait plus sous la phrase pensée.
\end{enumerate}

\subsection{Les statuts, et comment les lire}\label{sec:statuts}

\begin{table}[htbp]
\centering\small
\caption{Les statuts que les outils rendent — et ce qu'ils veulent dire.}
\label{tab:statuts}
\begin{tabularx}{\textwidth}{@{}l Y@{}}
\toprule
\textbf{Ce que vous voyez} & \textbf{Ce que ça veut dire} \\
\midrule
CF1 — statut VERT & le conteneur CF1 est trouvé, l'en-tête est conforme
(TRIADE-CF1 v1, Argon2id) — l'étiquette seule, aucune clé dérivée \\
CF1 — verifier INTÈGRE & la sentinelle a déverrouillé à froid : AEAD, empreinte,
projets — le coffre est entier \\
\code{[CF3] VERDICT : VERT — époque N, checksum a1b2c3d4, 3 projet(s)} & l'état
exact de votre coffre CF3 : époque à suivre, checksum à recouper, compte de projets
à connaître \\
AVERTISSEMENT D120 : SANS sentinelle vivante & vous déverrouillez un CF2 sur une
plateforme sans moteur : l'ère intermédiaire, honnête — close par CF3 \\
modules REQUIS absents & vous n'êtes pas dans une maison : le relais ne trouve pas
ses vérités vivantes — retournez au § \ref{sec:maison} \\
ÉCHEC NOMMÉ — le média du piège est REQUIS & branchez la clef USB du sachet et
désignez-la : la machine ne le crée jamais \\
\bottomrule
\end{tabularx}
\end{table}

\subsection{Les rapports}\label{sec:rapports}

« Envoyer le rapport » écrit un \code{.txt} : l'écran (journal d'écran), la boîte
noire horodatée, les empreintes md5 des modules (la génération de code qui a tourné se
\emph{lit}), l'environnement (python, OS, pile). Aucun réseau, aucun secret, aucun
chemin de coffre réel. Le rapport est la matière première du constat terrain (D122) :
glissez-le dans la conversation, le laboratoire le consigne et pousse. L'export n'est
pas un prérequis de la cérémonie — la chaîne vit même si le rapport meurt — mais il
est la mémoire de votre geste.

\subsection{L'autotest — prouver la serrure sans toucher au coffre}\label{sec:autotest}

Le lanceur \code{lancer\_autotest\_cf3} (racine de la maison) exécute les
autocontrôles de la serrure CF3 sur des répertoires \emph{temporaires} : aucun coffre
réel ni phrase réelle n'est touché (discipline D105). C'est le geste de santé après
une mise à jour des modules — un autotest vert sur une génération d'empreintes
conforme, puis votre coffre vérifiable à froid par \code{verifier}. L'équivalent CLI
de référence :

\begin{lstlisting}
py -3 v06\coffre_cf3.py autotest
py -3 v06\coffre_cf3.py verifier --conteneur .keys\coffre_cf3.bin ^
        --piege-dir <média du piège> --phrase-stdin
py -3 v06\sceau_I.py verifier      (dix arbres VERTS, code 0)
\end{lstlisting}

Le GUI de migration et le GUI de cérémonie font exactement ces appels — la CLI est la
référence, le GUI le relais. Sur Windows, \code{py -3} précède l'interpréteur ; la
phrase arrive par STDIN (\code{--phrase-stdin}), jamais en argument.
